
\subsubsection{6.1 Key Findings and Their Interpretation}

\textbf{Finding 1: Any tagging creates a substantial, robust discoverability advantage.} The 22.6\% adoption advantage (IRR=1.2263) survives the primary platform-era confound with only 12.2\% attenuation. The single most impactful metadata action for a dataset creator is adding at least one keyword tag. The choice of tags matters less than the act of tagging, because the primary mechanism is binary search index membership.

\textbf{Finding 2: Tag count magnitude is associated with a genuine dose-response.} Within the tagged subset, IRR=1.5332 per log-unit with an attenuation ratio of 1.0007 establishes that additional tags are associated with more task registrations independently of platform era. Each additional tag represents an additional keyword query pathway, reaching different researcher subsets. This is the amplification arm of the threshold-plus-amplification structure.

\textbf{Finding 3: Comprehensive tagging (6+) is the dominant amplification tier.} The 6+ tag tier produces a 28.6\% adoption advantage versus reference, meaningfully above the 12.7--12.8\% for 1--5 tags. The bimodal tagging distribution (all-or-nothing) means that minimal tagging (1--2 tags) and moderate tagging (3--5 tags) produce effects that are statistically indistinguishable from each other. Repository interfaces that encourage minimal tag entry may not realize the full discoverability benefit; 6+ tags appears to be the practical amplification threshold.

\textbf{Finding 4: The mechanism survives extreme temporal collinearity.} The Cramér's V=0.823 collinearity between decade and has\_tags represents a severe potential confound --- nearly all 2010s datasets are tagged, nearly none of the 2020s datasets are. That the has\_tags effect survives with only 12.2\% attenuation, while the composite metadata score attenuated to non-significance (98.6\%) under identical controls, establishes that the binary tagging variable captures a within-decade structural property --- not a platform-era trend.

\subsubsection{6.2 Connection to Literature}

The binary threshold finding provides the first NB-2 IRR for keyword tagging on an ML repository, filling the quantification gap in Wilkinson et al. [2016]. The direction is consistent with Yang et al. [2024] and Lachmuth et al. [2025], though platform and metric differences preclude direct comparison. The mechanism verification via decade fixed effects (H-M1) provides a methodological template for observational platform studies: testing extreme temporal collinearity without abandoning fixed effect controls demonstrates that within-cohort variation can identify structural mechanisms even under severe between-cohort confounding.

\subsubsection{6.3 Limitations}

\textbf{L1: Cross-sectional design --- predictive, not causal.} Tags and task counts are observed simultaneously. The causal interpretation rests on OpenML's creator-only tagging architecture \citep{Vanschoren2014OpenML} as a structural temporal ordering argument, not verified by timestamps in the available corpus.

\textbf{L2: N\_tasks proxy measures registration breadth, not execution depth.} Task registration is a deliberate adoption signal --- it requires defining a target variable and task type --- but may diverge from execution-frequency adoption for some datasets.

\textbf{L3: Conditional adoption only --- hurdle component deferred.} The analysis is scoped to $N_{\text{tasks}}\geq 1$ (adoption intensity). The probability of any adoption at all is not estimated.

\textbf{L4: 12.2\% decade attenuation is partially uncontrolled.} The residual after fixed effects may include both the structural FAIR F1 mechanism and uncontrolled era confounding. The gate passed with large margin despite this.

\textbf{L5: H-M3 bin 1--2 sparsity (N=73) limits uniform categorical dose-response claim.} The low-count range is undetermined; the 6+ tier characterization is robust. The informative negative reflects a data distribution constraint, not a mechanism failure.

\textbf{L6: Reverse causality for continuous IV.} For the log tag count measure (H-M2 and H-M3), popular datasets attracting additional tags retroactively remains a possible alternative explanation. OpenML's creator-only tagging architecture partially mitigates this, but if community tag additions exist, they could introduce upward bias in log\_tag\_count\_p1 estimates. This is not testable without tag modification timestamps from the platform API.

\subsubsection{6.4 Broader Implications}

Evidence-grounded IRR estimates translate FAIR F1 compliance from aspirational to actionable. Repository designers can justify tag-requirement policies with quantitative adoption effects; dataset creators gain guidance grounded in platform-scale evidence. The NB-2 methodology with decade fixed effects and attenuation testing provides a replicable template for other repository studies.

Strategic tag farming, analogous to SEO manipulation, could emerge if tagging is widely understood to boost adoption metrics. Repository governance should monitor for tag quality degradation alongside quantity. Additionally, if discovery concentrates among heavily-tagged datasets, a Matthew effect may emerge --- more-discoverable datasets accumulating further engagement at the expense of high-quality untagged datasets that do not appear in keyword searches.

